\documentclass[a4paper,11pt]{letter}
\usepackage[utf8]{inputenc}
\usepackage[T1]{fontenc}
\usepackage{lmodern}
\usepackage[english]{babel}
\usepackage[colorlinks=true, urlcolor=blue]{hyperref}
\usepackage{geometry}
\geometry{left=1in, right=1in, top=1in, bottom=1in}

\name{Baha Khammari}
\address{D\"usseldorf, Germany \\
    \href{mailto:b.e.khammari@gmail.com}{b.e.khammari@gmail.com} \\
    +49\,1520\,9016956}

\begin{document}

\begin{letter}{%
    Prof.\ Dr.\ Ronald Bachmann \\
    Labor Markets, Education, Population \\
    RWI -- Leibniz Institute for Economic Research \\
    Hohenzollernstra\ss e 1--3, 45128 Essen
}

\opening{Dear Professor Bachmann,}

I am writing to apply for a doctoral researcher position in the ``Labor Markets, Education, Population'' department at RWI. I am completing my M.Sc.\ in Economics at the University of Cologne (expected September 2026, current GPA: 1.6/1.0) and will be available from October 2026. My research interests are in applied microeconomics, causal inference, and education/labor economics, with a methodological focus on difference-in-differences estimation and heterogeneous treatment effects.

My doctoral project asks which higher education access mechanism---scholarships, loans, or quotas---is most effective at reducing labor market inequality, and for whom. I use Brazil as an empirical setting because it implemented all three at national scale within a fifteen-year window, but the research questions and comparative framework transfer to other country settings, including Germany.

My master's thesis (grade: 1.0/1.0, supervised by Anna Person, M.Sc.) provides the first causal analysis of Brazil's ProUni scholarship program on formal-sector wages at the microregion level. I constructed a panel of 558 regions over 2005--2019 from administrative microdata ($\sim$5\,million RAIS observations) and applied the Callaway \& Sant'Anna (2021) heterogeneity-robust staggered DiD estimator with doubly-robust inference. The key finding---a non-monotonic dose--response with significant wage gains in low-exposure regions that vanish at higher intensities---was entirely masked by conventional two-way fixed effects. The project required managing large-scale administrative datasets via SQL/BigQuery and implementing recent econometric methods in R.

RWI's ``Labor Markets, Education, Population'' department is a strong fit for this work. Prof.\ Westphal's research on causal returns to education using linked survey and administrative data and Dr.\ Tamm's evaluations of education reforms and their labor market effects both connect to the core of my project. The RWI-UNI-SUBJECTS dataset, which links study programs across German universities with social security records, would also be a valuable resource for extending the comparative framework to the German context---a direction I am interested in pursuing alongside the Brazilian analysis.

The RGS Econ's structured coursework in microeconometrics and the opportunity to work within RWI's applied policy research environment would let me combine methodological depth with direct policy relevance. Before the thesis, I worked for two years at PwC Germany (Assurance Solutions), where I built variance analyses on large financial datasets. An Erasmus+ semester at RSM Erasmus University (Rotterdam) added coursework in asset pricing and derivatives valuation.

My academic CV, a research proposal, and a thesis abstract are enclosed. Anna Person, M.Sc.\ (master's thesis supervisor), Univ.-Prof.\ Dr.\ Andreas Schabert (bachelor's thesis supervisor), and Prof.\ Dr.\ Michael Krause (seminar instructor) have agreed to serve as referees; their contact details are in my CV.

Thank you for considering my application.

\closing{Yours sincerely,}

\end{letter}

\end{document}
